\documentclass[12pt]{article}
\usepackage{amsmath, amssymb}
\usepackage[margin=1in]{geometry}
\setlength{\parskip}{6pt}
\title{QUBO Formulation: LLM Cascade Routing}
\date{}
\begin{document}
\maketitle

\subsection*{LLM Cascade Routing}

\paragraph{Problem Statement.}
Given a set of user requests $\mathcal{R}$ and a cascade of available language models $\mathcal{M}$, the goal is to assign each request to exactly one model tier. Each assignment $(r, m)$ incurs an inference cost $C_{r,m}$ and yields a predicted quality score $Q_m$. Each request $r$ has a minimum quality requirement $T_r$. The objective is to minimize the total inference cost while ensuring that every assigned model satisfies the quality threshold of its request.

\paragraph{Decision Variables.}
We introduce binary variables $x_{r,m} \in \{0,1\}$ for all $r \in \mathcal{R}$ and $m \in \mathcal{M}$, where $x_{r,m} = 1$ if and only if request $r$ is assigned to model $m$, and $0$ otherwise.

\paragraph{QUBO Derivation.}
We construct a Quadratic Unconstrained Binary Optimization (QUBO) formulation by incorporating quality constraints into the objective function via large penalties and enforcing the routing constraint via a squared penalty term.

Let $\tilde{C}_{r,m}$ denote the effective cost for assignment $(r,m)$, defined as:
\[
\tilde{C}_{r,m} = C_{r,m} + P \cdot \mathbb{I}(Q_m < T_r),
\]
where $P$ is a sufficiently large penalty constant and $\mathbb{I}(\cdot)$ is the indicator function. This ensures that infeasible assignments (those failing the quality threshold) carry a prohibitively high cost.

\textbf{Step 1: Objective Function.}
The minimization objective is:
\[
\min \sum_{r \in \mathcal{R}} \sum_{m \in \mathcal{M}} \tilde{C}_{r,m} x_{r,m}.
\]

\textbf{Step 2: Constraints.}
Each request must be routed to exactly one model:
\[
\sum_{m \in \mathcal{M}} x_{r,m} = 1, \quad \forall r \in \mathcal{R}.
\]

\textbf{Step 3: Penalty Terms.}
We enforce the equality constraints using a penalty weight $\lambda > 0$. For each request $r$, the penalty term is:
\[
H_r^{\text{one-hot}} = \lambda \left( \sum_{m \in \mathcal{M}} x_{r,m} - 1 \right)^2.
\]
Expanding the square and using the idempotent property $x_{r,m}^2 = x_{r,m}$ for binary variables:
\begin{align}
H_r^{\text{one-hot}} &= \lambda \left( \sum_{m \in \mathcal{M}} x_{r,m}^2 - 2 \sum_{m \in \mathcal{M}} x_{r,m} + 1 \right) \\
&= \lambda \left( \sum_{m \in \mathcal{M}} x_{r,m} - 2 \sum_{m \in \mathcal{M}} x_{r,m} + 1 \right) \\
&= \lambda \left( 1 - \sum_{m \in \mathcal{M}} x_{r,m} \right).
\end{align}
Alternatively, retaining the quadratic form for explicit QUBO structure:
\begin{align}
H_r^{\text{one-hot}} &= \lambda \left( 1 - \sum_{m \in \mathcal{M}} x_{r,m} + 2 \sum_{m < m' \in \mathcal{M}} x_{r,m} x_{r,m'} \right).
\end{align}

\textbf{Step 4: Hamiltonian.}
The total QUBO Hamiltonian $H(x)$ combines the objective and all penalty terms:
\begin{align}
H(x) &= \sum_{r \in \mathcal{R}} \sum_{m \in \mathcal{M}} \tilde{C}_{r,m} x_{r,m} + \sum_{r \in \mathcal{R}} \lambda \left( 1 - \sum_{m \in \mathcal{M}} x_{r,m} + 2 \sum_{m < m' \in \mathcal{M}} x_{r,m} x_{r,m'} \right).
\end{align}

\textbf{Step 5: Q Matrix and Offset.}
Collecting terms, the QUBO parameters are given by:
\begin{align}
Q_{(r,m),(r,m)} &= \tilde{C}_{r,m} - \lambda, \quad \forall r \in \mathcal{R}, m \in \mathcal{M}, \\
Q_{(r,m),(r,m')} &= 2\lambda, \quad \forall r \in \mathcal{R}, m, m' \in \mathcal{M}, m \neq m', \\
Q_{(r,m),(r',m')} &= 0, \quad \forall r, r' \in \mathcal{R}, r \neq r', \text{ or } m=m'.
\end{align}
The constant offset is:
\[
c = \sum_{r \in \mathcal{R}} \lambda = |\mathcal{R}| \lambda.
\]

\paragraph{Penalty Weight Justification.}
The penalty weight $\lambda$ must be chosen such that any violation of the routing constraint is strictly more costly than the maximum possible reduction in the objective function. Let $W = \max_{r,m} |\tilde{C}_{r,m}|$. If a solution violates a constraint, the penalty incurred is at least $\lambda$. The maximum gain in the objective from such a violation is bounded by $W$. Therefore, setting $\lambda > W$ ensures that the energy of any infeasible solution exceeds that of any feasible solution with non-negative cost, guaranteeing that the global minimum corresponds to a feasible assignment.

\paragraph{Correctness Argument.}
Minimizing $H(x)$ over $\{0,1\}^{|\mathcal{R}||\mathcal{M}|}$ recovers the optimal solution because: (i) for any feasible assignment, the penalty terms vanish, and $H(x)$ equals the total effective cost; (ii) for any infeasible assignment, at least one constraint is violated, incurring a penalty of at least $\lambda$. Since $\lambda > W$, the energy of any infeasible solution is strictly greater than the maximum possible cost of a feasible solution (assuming non-negative costs). Thus, the global minimum of $H(x)$ must be feasible and minimize the total cost.

\paragraph{Illustrative Example.}
Consider an instance with $|\mathcal{R}|=2$ requests and $|\mathcal{M}|=2$ models. Let the costs be $C_{0,0}=10$, $C_{0,1}=20$, $C_{1,0}=15$, $C_{1,1}=5$. Let quality scores be $Q_0=5$, $Q_1=8$, with thresholds $T_0=6$, $T_1=4$. The assignment $(0,0)$ is infeasible since $Q_0 < T_0$. Setting $P=1000$, we have $\tilde{C}_{0,0}=1010$ and $\tilde{C}_{r,m}=C_{r,m}$ otherwise. Let $\lambda=1000$.

The effective costs are $\tilde{C}_{0,0}=1010$, $\tilde{C}_{0,1}=20$, $\tilde{C}_{1,0}=15$, $\tilde{C}_{1,1}=5$. The Q matrix entries are:
\begin{align}
Q_{(0,0),(0,0)} &= 1010 - 1000 = 10, \\
Q_{(0,1),(0,1)} &= 20 - 1000 = -980, \\
Q_{(1,0),(1,0)} &= 15 - 1000 = -985, \\
Q_{(1,1),(1,1)} &= 5 - 1000 = -995, \\
Q_{(0,0),(0,1)} &= 2000, \quad Q_{(1,0),(1,1)} = 2000,
\end{align}
with all other off-diagonal entries zero. The offset is $c = 2 \times 1000 = 2000$.

The Hamiltonian is:
\begin{align}
H(x) &= 10 x_{0,0} - 980 x_{0,1} - 985 x_{1,0} - 995 x_{1,1} + 2000 x_{0,0} x_{0,1} + 2000 x_{1,0} x_{1,1} + 2000.
\end{align}
The optimal assignment is $x_{0,1}=1$, $x_{1,1}=1$ (requests routed to models 1 and 1), yielding energy $H = -980 - 995 + 2000 = 25$, which matches the total cost $20+5=25$.

\end{document}